\documentclass[italian,12pt]{article}

\def\Title{Verbale interno}
\def\Subtitle{Riunione interna settimanale}
\def\Author{7Last}
\def\Date{2024-05-22}
\def\Version{v1.0}

%pacchetti extra da scaricare dblfloatfix, fancyhdr
\usepackage[left=2cm, right=2cm, bottom=3cm, top=3cm]{geometry}
\usepackage{fancyhdr}%creazione header-footer
\usepackage{graphicx} %serve per inserire immagini
\graphicspath{ {../../../logo/} }
%\usepackage{dblfloatfix} %serve per posizionare gli elementi dove si vuole
\usepackage[hidelinks]{hyperref} %serve per i link
\usepackage{tikz}
\usepackage{tgadventor} % font
\usepackage[useregional=numeric,showseconds=true,showzone=false]{datetime2}
\usepackage{caption}

\usepackage{hyperref}
\usepackage{tocloft}
\usepackage{titlesec}
\usepackage{color}
\usepackage{ulem}
\usepackage{pgfplots}
\usepackage{pgf-pie}
\usepackage[italian]{babel}
\usepackage{comment}
\usepackage{tabularx}
\usepackage{longtable}
\usepackage{float}
\usepackage{amsmath}

\linespread{1.2}
\captionsetup[table]{labelformat=empty}
\geometry{headsep=1.5cm}

\renewcommand{\contentsname}{Indice}
\renewcommand\familydefault{\sfdefault}

\let\oldthepage\thepage
\renewcommand{\thepage}{\sffamily\oldthepage}


\begin{document}

\newgeometry{left=2cm,right=2cm,bottom=2.1cm,top=2.1cm}
\begin{titlepage}
    \vspace*{.5cm}

    \vspace{2cm}
    {
        \centering
        {\bfseries\huge \Title\par}
        \bigbreak
        {\bfseries\large \Author\par}
        \bigbreak
        {\Version\par}
        \vfill

        \begin{tikzpicture}[remember picture,overlay]

            \fill[blue!20!black] (current page.south west) -- ++(10cm,0) -- ++(-10cm,15cm) -- cycle;
            \fill[orange] (current page.south east) -- ++(-18cm,0) -- ++(21.6cm,22cm) -- cycle;

            \clip (0,-3cm) circle (2.5cm) node (current page.center) {\includegraphics[width=5cm]{logo.jpg}};
        \end{tikzpicture}

    }

    \vfill

\end{titlepage}

\restoregeometry






















\newpage

%---------------header------------------%
\pagestyle{fancy}
\fancyhead{} % pulizia degli header
\lhead{%
\begin{tikzpicture}
    \clip (0,0) circle (0.5cm);
    \node at (0,0) {\includegraphics[width=1cm]{logo}};
\end{tikzpicture}%
}
\chead{\vspace{\fill}\Title\vspace{\fill}}
\rhead{\vspace{\fill}\Version\vspace{\fill}}
%quad è una spaziatura
%---------------------------------------%

\begin{table}[!h]
	\caption{Versioni}
	\footnotesize
	\begin{center}
		\begin{tabular}{ l l l l p{6cm} }
			\hline                                                                           \\[-2ex]
			Ver. & Data       & Redattore      & Verificatore & Descrizione     \\
			\\[-2ex] \hline \\[-1.5ex]
			1.0  & 2024-06-09 & Leonardo Baldo & Elena Ferro  & Stesura verbale \\
			\\[-1.5ex] \hline
		\end{tabular}
	\end{center}
\end{table}
\newpage

\tableofcontents

\newpage

\section{Dettagli della riunione}


\textbf{Sede della riunione}: Piattaforma Discord\\
\textbf{Orario di inizio}: 17:00\\
\textbf{Orario di fine}: 18:00\\


\begin{flushleft}
	\begin{table}[!h]
		\begin{tabular}{ |l|l|l| }
			\hline
			\textbf{Partecipante} & \textbf{Ruolo} & \textbf{Presenza} \\
			\hline
			Antonio Benetazzo     &                & Presente          \\
			Davide Malgarise      & Amministratore & Presente          \\
			Elena Ferro           & Verificatore   & Presente          \\
			Leonardo Baldo        & Redattore      & Presente          \\
			Matteo Tiozzo         &                & Presente          \\
			Raul Seganfreddo      &                & Presente          \\
			Valerio Occhinegro    &                & Presente          \\
			\hline
		\end{tabular}
	\end{table}
\end{flushleft}

\section{Ordine del giorno}
\begin{itemize}
	\item Suddivisione dei ruoli per il settimo sprint
	\item Cosa è stato fatto fino ad ora
	\item Cosa rimane da fare
	\item Varie ed eventuali
	\item Decisioni prese e conclusioni
\end{itemize}

\newpage

\section{Verbale}
\subsection{Suddivisione dei ruoli per il settimo sprint}
Il gruppo stabilisce innanzitutto la suddivisione dei ruoli per il settimo periodo, assegnando a ciascun membro un compito non ancora svolto. La suddivisione risultante è la seguente:
\begin{itemize}
	\item \textbf{Responsabile}: Leonardo Baldo;
	\item \textbf{Amministratore}: Davide Malgarise;
	\item \textbf{Analista}: Raul Seganfreddo;
	\item \textbf{Progettista}: Matteo Tiozzo, Davide Malgarise, Valerio Occhinegro;
	\item \textbf{Programmatore}: Antonio Benetazzo;
	\item \textbf{Verificatore}: Elena Ferro.
\end{itemize}
Successivamente vengono definiti i ruoli per ogni riunione interna o esterna:
\begin{itemize}
	\item \textbf{Amministratore}: Davide Malgarise;
	\item \textbf{Redattore}: Leonardo Baldo;
	\item \textbf{Verificatore}: Elena Ferro.
\end{itemize}

\subsection{Cosa è stato fatto fino ad ora}
Viene fatta una panoramica sul lavoro svolto fino ad ora:
\begin{itemize}
	\item \textbf{Verbali} \\
		  Redatto, revisionato e approvato il verbale della riunione interna del 2024-05-29.
	\item \textbf{Stato della documentazione} \\
		  Continua la stesura dei documenti \textit{Piano di Qualifica}, aggiornando i grafici con i dati dello sprint e \textit{Piano di Progetto}, aggiungendo il consuntivo.
	\item \textbf{Studio} \\
		  Inoltre come prossimo passo, inizia lo studio dell'architettura di progetto.
\end{itemize}

\subsection{Cosa rimane da fare}
Per questo sprint ci sono ancora alcune attività da completare, riportate di seguito:
\begin{itemize}
	\item \textbf{Verbali} \\
		  Stesura del verbale della riunione in corso, ovvero della riunione interna del 2024-06-05.
	\item \textbf{Documenti} \\
		  Compilare il \textit{Piano di Progetto} per questo sprint e il \textit{Piano di Qualifica}, aggiungendo il consuntivo e i grafici con i dati dello sprint. \\
		  Inoltre inizia la stesura dei prossimi documenti, ovvero quelli relativi al \textit{Manuale Utente} e alla \textit{Scheda Tecnica}.
	\item \textbf{Codice} \\
		  Come richiesto dal proponente nello scorso SAL del 2024-06-05, viene chiesto di aggiornatre il codice del progetto, aggiungendo alcune tipologie di sensori nuovi, ovvero un sensore per l'umidità, uno per il parcheggio e uno per la qualità dell'aria. \\
		  Inoltre viene chiesto di iniziare a implementare Unit Test per il simulatore, per arrivare ad avere un Code Coverage.
\end{itemize}

\subsection{Varie ed eventuali}
Sono sorti alcuni dubbi e difficoltà riguardanti lo studio architetturale. Soprattutto nella individuazione di una architettura di progetto che meglio si adatti al nostro scopo. \\
Procediamo a continuarne lo studio per analizzare il problema e ad individuare una soluzione.

\subsection{Decisioni prese e conclusioni}
Non essendo emerse ulteriori problematiche o dubbi, il gruppo concorda la prossima riunione per il 2024-06-12.
\end{document}